\subsection{An exact linear obstruction at the inverse limit}
\label{sec:limit-relation-module}

The remaining intersection can be expressed by the actual
torsion modules and their quotient maps. Put
\[
 K=\bigcap_{n\geq1}\mathcal F_n=\ker\rho_{\mathrm{outer}},\qquad
 T_n=(\mathcal R_n)_{\Lambda\text{-tor}}
       =\mathcal F_{n+1}/J_{n+1},\qquad
 \mathcal L=\varprojlim_nT_n.
\]
The transition \(T_{n+1}\to T_n\) is induced by inclusion
\(\mathcal F_{n+2}\subseteq \mathcal F_{n+1}\), followed by reduction modulo
\(J_{n+1}\).

\begin{proposition}\label{prop:limit-relation-module}
There are canonical topological isomorphisms
\begin{align}
 K/\overline{[H,K]}&\simeq\mathcal L,
 \label{eq:limit-relation-iso}\\
 \mathcal L/(p,\gamma-1)\mathcal L
 &\simeq
 \varprojlim_n\bigl(T_n/(p,\gamma-1)T_n\bigr).
 \label{eq:limit-mod-p-iso}
\end{align}
Both modules on the right detect the triviality of \(K\):
\begin{equation}\label{eq:limit-linear-obstruction}
 K=1\quad\Longleftrightarrow\quad
 \mathcal L=0\quad\Longleftrightarrow\quad
 \varprojlim_n\bigl(T_n/(p,\gamma-1)T_n\bigr)=0.
\end{equation}
Writing \(A_n=\operatorname{im}(T_n\to M)\) and
\(P_n=\ker(T_n\to M)\), there is also an exact sequence
\begin{equation}\label{eq:limit-abelian-torsion}
 0\longrightarrow\varprojlim_nP_n
 \longrightarrow\mathcal L
 \longrightarrow\bigcap_nA_n
 \longrightarrow0,
 \qquad p^{h_n}P_n=0.
\end{equation}
\end{proposition}

\begin{proof}
For \(n\geq2\), the commutator definitions give
\[
 \overline{[H,K]}\subseteq J_n
 \subseteq\overline{[H,\mathcal F_{\lceil n/2\rceil}]}.
\]
In any finite quotient of \(H\), the descending images of
\(\mathcal F_m\) eventually stabilize at the image of \(K\).
Indeed every coset belonging to all these images has
nonempty nested compact intersections with the \(\mathcal F_m\).
Consequently the images of the two commutator bounds
eventually coincide. Passing through all finite quotients
proves
\[
                    \bigcap_nJ_n=\overline{[H,K]}.
\]
A compatible family in \(\varprojlim T_n\) is a nested family
of nonempty compact cosets of \(J_{n+1}\), each contained
in \(\mathcal F_{n+1}\). A point of their intersection belongs to
\(K\). The kernel of the resulting surjection
\(K\to\mathcal L\) is the displayed intersection of the
\(J_n\). Compactness proves
\eqref{eq:limit-relation-iso} as a topological isomorphism.

If \(\mathcal L=0\), then \(K=\overline{[H,K]}\).
Its image \(N\) in any finite \(p\)-group quotient \(G\)
of \(H\) satisfies \(N=[G,N]\). Iteration puts \(N\)
in every lower central term of the nilpotent group \(G\),
so \(N=1\). Every finite image of \(K\) is trivial,
which proves \(K=1\).

Inverse limits of these exact sequences of compact groups
are exact. For the surjectivity assertion, a compatible
family of target points can be lifted over any finite
initial set of indices by lifting first at the largest
index and pushing down. The resulting finite systems
of closed conditions have the finite intersection
property in a product of compact groups.
Apply this observation to the maps
\[
 T_n^2\longrightarrow T_n,\qquad
 (x,y)\longmapsto px+(\gamma-1)y,
\]
their compact images, and their cokernels. It proves
\eqref{eq:limit-mod-p-iso}.

For completeness, compact Nakayama in this setting does
not need finite generation. Open \(\Gamma\)-stable
subgroups form a neighborhood basis in \(\mathcal L\):
an open subgroup has an open stabilizer in the compact
group \(\Gamma\), and one intersects its finitely many
translates. On a finite quotient by such a subgroup,
\((p,\gamma-1)\) acts nilpotently, since \(\Gamma\)
has finite \(p\)-group image. Thus
\(\mathcal L=(p,\gamma-1)\mathcal L\) makes every
finite quotient zero. This proves
\eqref{eq:limit-linear-obstruction}.

Finally the maps \(T_n\to M\) have nested images \(A_n\).
The compact exact sequences
\(0\to P_n\to T_n\to A_n\to0\) give
\eqref{eq:limit-abelian-torsion}. The bound on \(P_n\)
is the finite-stage kernel bound already proved.
Equivalently, nested compact fibers give directly
\(\bigcap A_n=\operatorname{im}(K\to M)\).
\end{proof}

The exponents \(h_n\) are finite at each stage; no uniform
bound in \(n\) is asserted. Thus
\eqref{eq:limit-abelian-torsion} retains the inverse limit
of the inherited bounded torsion as well as the intersection
of the finite-rank lattices. The last module in
\eqref{eq:limit-linear-obstruction}, with its actual
transition maps, is the precise remaining linear
vanishing question.

\paragraph{The homology measured by the torsion limit.}
Put \(G=H/K=\rho_{\mathrm{outer}}(H)\), with its actual compact
quotient topology. The free pro-\(p\) presentation
\(1\to K\to H\to G\to1\) gives the continuous Hopf sequence
\[
0\longrightarrow H_2^{\mathrm{cont}}(G,\mathbb Z_p)
\longrightarrow K/\overline{[H,K]}
\longrightarrow H^{\mathrm{ab}}
\longrightarrow G^{\mathrm{ab}}\longrightarrow0.
\]
For the precise pro-\(p\) morphism, see
\cite[Proposition~2.2, author p.~3, and the Hopf formula on
author p.~4]{Moravec2011}. It also follows directly by
Pontryagin-dualizing the Hochschild--Serre five-term sequence
with discrete coefficients \(\mathbb Q_p/\mathbb Z_p\);
freeness of \(H\) makes its second cohomology zero.
This argument applies to the present countable source and
does not require a finite presentation of \(G\).
The kernel of the middle map is
\((K\cap\overline{[H,H]})/\overline{[H,K]}\).
Comparison with \eqref{eq:limit-abelian-torsion} therefore gives
the actual topological isomorphism
\[
 \varprojlim_nP_n
       \simeq H_2^{\mathrm{cont}}(G,\mathbb Z_p).
\]
The finite-stage bounds on \(P_n\) thus identify an exact
continuous homology term; they do not evaluate its inverse
limit.

